%===============================================================================
%  Code Review Report: Enlsip.jl
%  Analysis against Fortran 77 ENLSIP and Mathematical Specification
%===============================================================================
\documentclass[11pt, a4paper]{article}

\usepackage{amsmath, amssymb}
\usepackage{geometry}
\usepackage{hyperref}
\usepackage{booktabs}
\usepackage{array}
\usepackage{xcolor}
\usepackage{listings}
\usepackage{enumitem}
\usepackage{fancyhdr}
\usepackage{caption}

\geometry{margin=2.5cm}

\hypersetup{
  colorlinks=true,
  linkcolor=blue!70!black,
  citecolor=green!50!black,
  urlcolor=blue!70!black,
}

\lstdefinelanguage{Julia}{
  morekeywords={abstract,baremodule,begin,bitstype,break,catch,ccall,const,
    continue,do,else,elseif,end,export,finally,for,function,global,if,
    immutable,import,importall,in,let,local,macro,module,mutable,
    nothing,quote,return,struct,try,type,typealias,using,where,while,
    true,false},
  sensitive=true,
  morecomment=[l]{\#},
  morecomment=[n]{\#=}{=\#},
  morestring=[b]",
}
\lstset{
  language=Julia,
  basicstyle=\ttfamily\small,
  keywordstyle=\color{blue!70!black}\bfseries,
  commentstyle=\color{green!50!black}\itshape,
  stringstyle=\color{red!60!black},
  showstringspaces=false,
  breaklines=true,
  frame=single,
  framerule=0.4pt,
  rulecolor=\color{gray!50},
  backgroundcolor=\color{gray!5},
  numbers=left,
  numberstyle=\tiny\color{gray},
  numbersep=6pt,
  tabsize=4,
  columns=flexible,
}

\newcommand{\fortranref}[1]{\texttt{#1}}
\newcommand{\juliaref}[2]{\texttt{#1:#2}}
\newcommand{\severity}[1]{%
  \ifx#1C\colorbox{red!20}{\textbf{Critical}}%
  \else\ifx#1M\colorbox{orange!20}{\textbf{Major}}%
  \else\ifx#1m\colorbox{yellow!20}{\textbf{Minor}}%
  \else\colorbox{blue!10}{\textbf{Enhancement}}%
  \fi\fi\fi%
}

\pagestyle{fancy}
\fancyhf{}
\rhead{\small Code Review: Enlsip.jl}
\lhead{\small \today}
\rfoot{\thepage}

\title{%
  \textbf{Code Review Report: Enlsip.jl}\\[8pt]
  \large Cross-Reference Analysis against Fortran~77 ENLSIP\\
  and Algorithmic Specification
}
\author{Automated review based on comparison with:\\[4pt]
  \textit{Fortran~77 ENLSIP} (Lindstr\"om \& Wedin, 1988)\\
  \textit{ENLSIP.jl} (reference Julia translation)\\
  \textit{Mathematical Analysis of ENLSIP} (convergence analysis document)
}
\date{\today}

\begin{document}
\maketitle

\begin{abstract}
This report presents the results of a detailed code review of the
\texttt{Enlsip.jl} Julia package, a published implementation of the
ENLSIP constrained nonlinear least-squares solver.  The review was
conducted by cross-referencing the Julia source against the original
Fortran~77 implementation, a reference Julia translation, and a
mathematical analysis of the algorithm.  We identify 7~bugs (3~critical,
4~major), 5~algorithmic discrepancies from the original, and
6~improvement opportunities.  For each finding, we describe the issue,
its impact, the relevant Fortran reference, and the applied correction.
\end{abstract}

\tableofcontents
\newpage

%===============================================================================
\section{Overview}
%===============================================================================

The \texttt{Enlsip.jl} package is a Julia implementation of the ENLSIP
algorithm for solving constrained nonlinear least-squares problems of the
form:
\begin{align}
  \min_{\mathbf{x} \in \mathbb{R}^n} \quad & \Phi(\mathbf{x}) = \tfrac{1}{2}\|\mathbf{r}(\mathbf{x})\|^2 \\
  \text{subject to} \quad & h_i(\mathbf{x}) = 0, \quad i = 1,\dots,p \\
  & h_j(\mathbf{x}) \geq 0, \quad j = p+1,\dots,\ell
\end{align}
where $\mathbf{r} : \mathbb{R}^n \to \mathbb{R}^m$ and
$\mathbf{h} : \mathbb{R}^n \to \mathbb{R}^\ell$, with $m + p \geq n$.

The review covers the four source files:
\begin{itemize}[nosep]
  \item \texttt{src/structures.jl} --- data structures (\texttt{Iteration}, \texttt{WorkingSet}, etc.)
  \item \texttt{src/cnls\_model.jl} --- problem model, evaluation wrappers, constraint instantiation
  \item \texttt{src/enlsip\_functions.jl} --- core algorithmic routines ($\sim$2750~lines)
  \item \texttt{src/solver.jl} --- user-facing \texttt{solve!} interface
\end{itemize}

Findings are classified as:
\begin{itemize}[nosep]
  \item \severity{C} --- silent incorrect results or crash on valid input
  \item \severity{M} --- incorrect behaviour in non-trivial edge cases
  \item \severity{m} --- cosmetic or minor correctness issues
  \item \severity{E} --- performance or design improvements
\end{itemize}

%===============================================================================
\section{Critical Bugs}
\label{sec:critical}
%===============================================================================

%-----------------------------------------------------------------------
\subsection{B1: Integer division truncation in \texttt{hessian\_cons!}}
\label{sec:b1}
%-----------------------------------------------------------------------

\severity{C} \quad \juliaref{enlsip\_functions.jl}{296}
\quad Fortran ref: \fortranref{HESSH}

\paragraph{Issue.}
The finite-difference step size for the constraint Hessian approximation
is computed as:
\begin{lstlisting}
  eps1 = eps(T)^(1 / 3)
\end{lstlisting}
In Julia, \texttt{1 / 3} performs \emph{integer} division and evaluates
to~\texttt{0}.  Therefore \texttt{eps1 = eps(T)\^{}0 = 1.0}, which is
entirely wrong --- the step size should be $\varepsilon_{\mathrm{mach}}^{1/3}
\approx 6 \times 10^{-6}$.

\paragraph{Impact.}
Every Newton search direction (method code~2) that involves the
constraint Hessian will use a step size of~1.0 for finite differencing,
producing a completely inaccurate second-derivative approximation.  This
silently corrupts Newton steps and can cause convergence failure or
convergence to wrong points on problems requiring second-order
corrections.

\paragraph{Comparison.}
The analogous function \texttt{hessian\_res!} at line~252 correctly uses
\texttt{eps(T)\^{}(1.0 / 3.0)} with floating-point literals.  The
Fortran code uses \texttt{DRELPR**(1.0D0/3.0D0)}.

\paragraph{Fix.}
\begin{lstlisting}
  eps1 = eps(T)^(1.0 / 3.0)
\end{lstlisting}

%-----------------------------------------------------------------------
\subsection{B2: Undefined variable names in constraint Jacobian instantiation}
\label{sec:b2}
%-----------------------------------------------------------------------

\severity{C} \quad \juliaref{cnls\_model.jl}{488, 492, 439}

\paragraph{Issue.}
In \texttt{instantiate\_constraints\_wo\_bounds}, two code paths define a
local Jacobian function but then pass an \emph{undefined} variable name
to the constructor:

\begin{lstlisting}
  # Line 488: defines jac_cons_ineqnum, passes jac_ineqnum (undefined)
  jac_cons_ineqnum(x) = vcat(jacobian_eqcons(x),
      ForwardDiff.jacobian(ineq_constraints, x))
  constraints_evalfunc = ConstraintsFunction(cons, jac_ineqnum)

  # Line 492: defines ad_jac_eqcons, passes jac_eqnum (undefined)
  ad_jac_eqcons(x) = vcat(ForwardDiff.jacobian(eq_constraints,x),
      jacobian_ineqcons(x))
  constraints_evalfunc = ConstraintsFunction(cons, jac_eqnum)
\end{lstlisting}

The same pattern occurs in \texttt{instantiate\_constraints\_w\_bounds}
at line~439.

\paragraph{Impact.}
Any problem that provides exactly one of the two Jacobians (equality
\emph{xor} inequality) while omitting the other will trigger an
\texttt{UndefVarError} at runtime.  The bug is latent because all test
problems provide either both Jacobians or neither.

\paragraph{Fix.}
Replace each undefined name with the corresponding local variable that
was actually defined.

%-----------------------------------------------------------------------
\subsection{B3: Shallow copy of mutable vectors in \texttt{Iteration}}
\label{sec:b3}
%-----------------------------------------------------------------------

\severity{C} \quad \juliaref{structures.jl}{94--95}

\paragraph{Issue.}
The \texttt{Base.copy} method for \texttt{Iteration} performs a
field-by-field copy:
\begin{lstlisting}
  Base.copy(s::Iteration) = Iteration(s.x, s.p, s.rx, ...)
\end{lstlisting}
Since \texttt{s.x}, \texttt{s.p}, \texttt{s.rx}, etc.\ are mutable
\texttt{Vector}s, the ``copy'' shares the same underlying arrays.
Subsequent mutation of one iteration's vectors silently corrupts the
other.

\paragraph{Impact.}
In the main loop, \texttt{previous\_iter = copy(iter)} is intended to
snapshot the current state.  Because vectors are shared,
\texttt{previous\_iter.rx}, \texttt{previous\_iter.p}, etc.\ can be
overwritten when \texttt{iter} is updated, leading to incorrect
convergence checks, search direction analysis, and penalty weight
updates that depend on the previous iteration's data.

\paragraph{Fix.}
Use \texttt{deepcopy} for all vector fields:
\begin{lstlisting}
  Base.copy(s::Iteration) = Iteration(
      copy(s.x), copy(s.p), copy(s.rx), copy(s.cx), s.t, s.alpha,
      s.index_alpha_upp, copy(s.lambda), copy(s.w), ...)
\end{lstlisting}

%===============================================================================
\section{Major Bugs}
\label{sec:major}
%===============================================================================

%-----------------------------------------------------------------------
\subsection{B4: \texttt{second\_lagrange\_mult\_estimate!} ignores pseudo-rank}
\label{sec:b4}
%-----------------------------------------------------------------------

\severity{M} \quad \juliaref{enlsip\_functions.jl}{522}
\quad Fortran ref: \fortranref{LEAEST}

\paragraph{Issue.}
The second-order Lagrange multiplier estimate solves:
\begin{lstlisting}
  v = UpperTriangular(F_A.R) \ b
\end{lstlisting}
This uses the full $R$ factor regardless of the pseudo-rank of~$A$.
When $A$ is rank-deficient ($\mathrm{rank} < t$), the trailing diagonal
elements of~$R$ are near-zero, producing wildly inaccurate multiplier
estimates.

\paragraph{Comparison.}
The first-order estimate (\texttt{first\_lagrange\_mult\_estimate!},
lines~476--478) correctly restricts the solve to the leading
\texttt{prankA} $\times$ \texttt{prankA} block.

\paragraph{Fix.}
Compute pseudo-rank and restrict the triangular solve accordingly.

%-----------------------------------------------------------------------
\subsection{B5: \texttt{evaluate\_violated\_constraints} lacks capacity bound}
\label{sec:b5}
%-----------------------------------------------------------------------

\severity{M} \quad \juliaref{enlsip\_functions.jl}{600--623}
\quad Fortran ref: \fortranref{EVADD}

\paragraph{Issue.}
The Fortran \fortranref{EVADD} enforces the bound
$|\mathcal{T}| \leq \min(\ell, n)$ on the working-set size.  When the
working set is at capacity and a violated constraint must be added,
EVADD swaps out the least-violated active constraint to make room.

The Julia implementation has no such capacity check --- it adds all
violated constraints unconditionally.  This can produce $|\mathcal{T}| > n$,
making the constraint Jacobian $A_{\mathcal{T}}$ have more rows than columns.
The resulting LQ factorisation and null-space step become algebraically
ill-defined.

\paragraph{Fix.}
Add a capacity check \texttt{W.t >= min(W.l, n)} before adding constraints,
with a swap-out mechanism for the least-violated active constraint when at capacity.

%-----------------------------------------------------------------------
\subsection{B6: \texttt{check\_termination\_criteria} uses stale \texttt{iter.x}}
\label{sec:b6}
%-----------------------------------------------------------------------

\severity{M} \quad \juliaref{enlsip\_functions.jl}{2425}

\paragraph{Issue.}
In the abnormal-termination branch, the step difference is computed as:
\begin{lstlisting}
  x_diff = norm(prev_iter.x - iter.x)
\end{lstlisting}
However, \texttt{iter.x} is only updated \emph{after} the termination
check (main loop, line~2732: \texttt{iter.x = x}).  Therefore
\texttt{iter.x} still holds the previous iteration's starting point,
and the computed \texttt{x\_diff} is meaningless.

The sufficient-conditions branch (line~2401) correctly uses the passed
parameter \texttt{x}.

\paragraph{Fix.}
Use the parameter \texttt{x} consistently in both branches.

%-----------------------------------------------------------------------
\subsection{B7: \texttt{newton\_raphson} can loop infinitely}
\label{sec:b7}
%-----------------------------------------------------------------------

\severity{M} \quad \juliaref{enlsip\_functions.jl}{1764}

\paragraph{Issue.}
The Newton--Raphson loop for polynomial root-finding has the condition:
\begin{lstlisting}
  while error > eps || newton_iter < 3
\end{lstlisting}
The \texttt{||} operator means the loop continues as long as
\emph{either} condition holds.  If the polynomial derivative
\texttt{dds($\alpha$)} is near zero, the step $h = -ds(\alpha)/c$
diverges, \texttt{error} never decreases, and the loop runs forever.

There is no iteration cap.

\paragraph{Fix.}
Add a maximum iteration bound:
\begin{lstlisting}
  while (error > eps || newton_iter < 3) && newton_iter < 50
\end{lstlisting}

%===============================================================================
\section{Algorithmic Discrepancies}
\label{sec:discrepancies}
%===============================================================================

%-----------------------------------------------------------------------
\subsection{D1: \texttt{pseudo\_rank} scaling factor}
\label{sec:d1}
%-----------------------------------------------------------------------

\juliaref{enlsip\_functions.jl}{17--31}
\quad Fortran ref: \fortranref{PSEUDORANKR1}

The student's implementation uses the tolerance
$\mathrm{tol} = |R_{11}| \cdot \sqrt{\ell} \cdot \varepsilon_{\mathrm{rank}}$,
whereas the Fortran uses $\mathrm{tol} = |R_{11}| \cdot \varepsilon_{\mathrm{rank}}$.
The additional $\sqrt{\ell}$ factor makes the rank determination more
conservative.  While not necessarily incorrect, this deviation can
produce different pseudo-ranks on marginally rank-deficient problems,
affecting the entire algorithm flow.

\medskip\noindent
\textit{No change applied --- this is a deliberate design choice documented here
for awareness.}

%-----------------------------------------------------------------------
\subsection{D2: Missing anti-cycling guard in \texttt{check\_constraint\_deletion}}
\label{sec:d2}
%-----------------------------------------------------------------------

\juliaref{enlsip\_functions.jl}{567--596}
\quad Fortran ref: \fortranref{SIGNCH}, lines 2322--2328

The Fortran \fortranref{SIGNCH} tracks iteration numbers in the
\texttt{ACTIVE} array to prevent a constraint from being deleted and
re-added within a small number of iterations (the ``ival'' guard).
This anti-cycling mechanism prevents pathological oscillation of the
working set.

The Julia version has no such safeguard.  Implementing it would require
extending the \texttt{WorkingSet} structure to track deletion history.

\medskip\noindent
\textit{Not corrected in this revision --- would require structural changes to
\texttt{WorkingSet}.  Documented as a known limitation.}

%-----------------------------------------------------------------------
\subsection{D3: Missing infeasibility detection in convergence code}
\label{sec:d3}
%-----------------------------------------------------------------------

\juliaref{enlsip\_functions.jl}{2373--2456}
\quad Fortran ref: \fortranref{TERCRI}, lines 816--817

The Fortran TERCRI constructs the exit code as $\texttt{EXIT} \times
\texttt{feas}$, where $\texttt{feas} = -1$ when any inactive constraint
is violated.  This turns a positive convergence code into a negative
``converged to infeasible point'' code, which is a crucial diagnostic.

The Julia version returns the positive convergence code regardless of
feasibility.

\paragraph{Fix.}
After computing the sufficient-condition code, check inactive constraint
feasibility and negate the code if infeasible.

%-----------------------------------------------------------------------
\subsection{D4: Missing full-rank prerequisite for criterion~8}
\label{sec:d4}
%-----------------------------------------------------------------------

\juliaref{enlsip\_functions.jl}{2378}
\quad Fortran ref: \fortranref{TERCRI}, lines 815--832

The Fortran adds a ``last step type'' digit (criterion~8) to the exit
code based on the history of the last few steps (GN vs.\ Newton vs.\
subspace, full rank, unit step length).  The Julia version does not
implement this refinement, which is used for diagnostic purposes.

\medskip\noindent
\textit{Not corrected --- purely diagnostic, no effect on convergence decisions.}

%-----------------------------------------------------------------------
\subsection{D5: \texttt{minmax\_lagrangian\_mult} initial \texttt{sigmin} value}
\label{sec:d5}
%-----------------------------------------------------------------------

\juliaref{enlsip\_functions.jl}{545}

\texttt{sigmin} is initialised to $10^6$ rather than to $+\infty$ or a
value derived from the problem data.  If all inequality multipliers
satisfy $\lambda_i \cdot \mathrm{row}_i > -\sqrt{\varepsilon}$, the
initial value persists and $\mathrm{sigmin} = 10^6$ is compared against
$\varepsilon_{\mathrm{rel}} \cdot \mathrm{factor}$ in the convergence
test.  While $10^6$ is ``large enough'' in practice, using \texttt{Inf}
would be semantically correct.

\paragraph{Fix.}
Initialise \texttt{sigmin = Inf}.

%===============================================================================
\section{Performance and Design Improvements}
\label{sec:improvements}
%===============================================================================

%-----------------------------------------------------------------------
\subsection{I1: Excessive allocations in Hessian routines}
\label{sec:i1}
%-----------------------------------------------------------------------

\juliaref{enlsip\_functions.jl}{243--324}

The functions \texttt{hessian\_res!} and \texttt{hessian\_cons!}
allocate four vectors of length~$m$ (resp.~$\ell$) and two unit vectors
of length~$n$ at every $(k,j)$ pair in a double loop over
$1 \leq j \leq k \leq n$.  This creates $O(n^2)$ allocations of
$O(m)$-sized arrays, generating significant garbage-collection pressure
for large problems.

\paragraph{Fix.}
Pre-allocate work arrays outside the loops and reuse them.  Replace
the boolean unit-vector construction \texttt{[i == k for i = 1:n]} with
in-place manipulation of a single pre-allocated vector.

%-----------------------------------------------------------------------
\subsection{I2: \texttt{psi} re-evaluates functions on every call}
\label{sec:i2}
%-----------------------------------------------------------------------

\juliaref{enlsip\_functions.jl}{1269--1304}

The merit function \texttt{psi} allocates new vectors and evaluates both
the residual and constraint functions from scratch on each invocation.
During the line search, it may be called $O(10)$ times.  While the
polynomial-based line search (\texttt{minrm}/\texttt{minrn}) reduces the
number of evaluations, the fallback Goldstein--Armijo step
(\texttt{goldstein\_armijo\_step}) calls \texttt{psi} in a halving loop.

\medskip\noindent
\textit{Not modified in this revision --- would require interface changes
to pass pre-allocated work arrays.}

%-----------------------------------------------------------------------
\subsection{I3: \texttt{f\_opt} display is one iteration behind}
\label{sec:i3}
%-----------------------------------------------------------------------

\juliaref{enlsip\_functions.jl}{2719}

In the main loop, the \texttt{DisplayedInfo} object pushed to the
iteration-detail list uses \texttt{f\_opt}, but this variable is only
updated at line~2737 (\emph{after} the push).  The displayed objective
function value therefore lags by one iteration.

\paragraph{Fix.}
Update \texttt{f\_opt} before constructing the \texttt{DisplayedInfo}.

%-----------------------------------------------------------------------
\subsection{I4: No support for unconstrained problems}
\label{sec:i4}
%-----------------------------------------------------------------------

\juliaref{cnls\_model.jl}{377}

The \texttt{CnlsModel} constructor asserts that at least one constraint
exists.  The original Fortran ENLSIP handles unconstrained problems
($\ell = 0$) gracefully.

\medskip\noindent
\textit{Not modified --- this is a design decision for the package scope.}

%===============================================================================
\section{Summary of Applied Corrections}
\label{sec:summary}
%===============================================================================

Table~\ref{tab:summary} summarises all findings and their resolution status.

\begin{table}[ht]
\centering
\caption{Summary of findings and corrections.}
\label{tab:summary}
\small
\begin{tabular}{@{}llllp{5.2cm}@{}}
\toprule
\textbf{ID} & \textbf{Severity} & \textbf{File} & \textbf{Line(s)} & \textbf{Status} \\
\midrule
B1 & Critical & \texttt{enlsip\_functions.jl} & 296     & \textbf{Fixed}: \texttt{1/3} $\to$ \texttt{1.0/3.0} \\
B2 & Critical & \texttt{cnls\_model.jl}       & 438--492 & \textbf{Fixed}: corrected variable names \\
B3 & Critical & \texttt{structures.jl}        & 94--95   & \textbf{Fixed}: deep copy of vectors \\
B4 & Major    & \texttt{enlsip\_functions.jl} & 522     & \textbf{Fixed}: pseudo-rank--aware solve \\
B5 & Major    & \texttt{enlsip\_functions.jl} & 600--623 & \textbf{Fixed}: capacity check added \\
B6 & Major    & \texttt{enlsip\_functions.jl} & 2425    & \textbf{Fixed}: use \texttt{x} parameter \\
B7 & Major    & \texttt{enlsip\_functions.jl} & 1764    & \textbf{Fixed}: iteration cap added \\
\midrule
D1 & ---      & \texttt{enlsip\_functions.jl} & 17--31  & Documented only \\
D2 & ---      & \texttt{enlsip\_functions.jl} & 567--596 & Documented only \\
D3 & Major    & \texttt{enlsip\_functions.jl} & 2399--2421 & \textbf{Fixed}: infeasibility negation \\
D5 & Minor    & \texttt{enlsip\_functions.jl} & 545     & \textbf{Fixed}: \texttt{sigmin = Inf} \\
\midrule
I1 & Enh.     & \texttt{enlsip\_functions.jl} & 243--324 & \textbf{Fixed}: pre-allocated work arrays \\
I3 & Minor    & \texttt{enlsip\_functions.jl} & 2719    & \textbf{Fixed}: \texttt{f\_opt} update order \\
\bottomrule
\end{tabular}
\end{table}

%===============================================================================
\section{Testing Recommendations}
%===============================================================================

The existing test suite covers four problems (HS65, Osborne-2, Chained
Rosenbrock, Chained Wood) but exercises only a subset of the code paths.
We recommend adding tests for:

\begin{enumerate}[nosep]
  \item \textbf{Mixed Jacobian provision}: provide only the equality
        Jacobian and omit the inequality Jacobian (and vice versa) to
        exercise the paths fixed in~B2.
  \item \textbf{Rank-deficient constraint Jacobians}: problems where
        active constraints are linearly dependent at some iterate, to
        test the pseudo-rank handling in B4.
  \item \textbf{Large-residual problems}: problems where the
        Gauss--Newton model is insufficient and Newton steps are needed,
        to exercise the corrected Hessian computation in~B1.
  \item \textbf{Working-set saturation}: problems with $\ell > n$
        inequality constraints where the capacity bound (B5) is reached.
  \item \textbf{Convergence to infeasible points}: crafted problems
        where the algorithm converges but inactive constraints are
        violated, to test the infeasibility detection in~D3.
\end{enumerate}

%===============================================================================
\section{Performance Optimisations}
\label{sec:optim}
%===============================================================================

This section describes performance-oriented changes applied in a second
pass.  Unlike the bug-fix pass, these changes do not alter algorithmic
behaviour but reduce memory allocations, improve type stability, and
eliminate redundant computation.

%-----------------------------------------------------------------------
\subsection{O1: Type-unstable function fields in \texttt{ResidualsFunction}
  and \texttt{ConstraintsFunction}}
\label{sec:o1}
%-----------------------------------------------------------------------

\juliaref{cnls\_model.jl}{11--35}

Both \texttt{ResidualsFunction} and \texttt{ConstraintsFunction} store
their evaluation and Jacobian callables as untyped (\texttt{Any}) fields.
Every call through these fields forces Julia to use dynamic dispatch,
which prevents inlining and generates type-instabilities that propagate
throughout the solver.

\paragraph{Fix.}
Parameterise the structs over two function types:
\begin{lstlisting}
  mutable struct ResidualsFunction{F,G} <: EvaluationFunction
      reseval::F
      jacres_eval::G
      nb_reseval::Int
      nb_jacres_eval::Int
  end
\end{lstlisting}
Since both are set once at construction and never change type,
this lets Julia specialise all downstream code.

%-----------------------------------------------------------------------
\subsection{O2: Untyped function fields in \texttt{CnlsModel}}
\label{sec:o2}
%-----------------------------------------------------------------------

\juliaref{cnls\_model.jl}{155--174}

The \texttt{CnlsModel} struct stores six function/nothing fields
(\texttt{residuals}, \texttt{jacobian\_residuals}, \texttt{eq\_constraints},
etc.) as untyped \texttt{Any} fields.  These fields are only read at
\texttt{solve!} time to construct the typed evaluation functions, so the
impact during the inner loop is indirect (the eval functions are already
wrapped).  Nonetheless, parameterising the model struct improves
type stability of the \texttt{solve!} entry point.

\medskip\noindent
\textit{Not modified --- the performance gain is negligible since these
fields are not accessed in the inner loop.  Changing the struct signature
would break the public API.}

%-----------------------------------------------------------------------
\subsection{O3: \texttt{map(abs, $\cdot$)} $\to$ \texttt{maximum(abs, $\cdot$)}}
\label{sec:o3}
%-----------------------------------------------------------------------

\juliaref{enlsip\_functions.jl}{554, 585, 1553, 2301}

The pattern \texttt{maximum(map(abs, v))} allocates an intermediate array;
\texttt{maximum(abs, v)} (Julia~$\geq$~1.6) computes the same thing
without allocation.  Similarly, \texttt{maximum(map(t -> abs(t), $\lambda$))}
should be \texttt{maximum(abs, $\lambda$)}.

\paragraph{Fix.}
Replace all four occurrences.

%-----------------------------------------------------------------------
\subsection{O4: \texttt{setdiff(1:end, s)} row-deletion pattern}
\label{sec:o4}
%-----------------------------------------------------------------------

\juliaref{enlsip\_functions.jl}{719, 752, 756, 780, 785}

Deleting row~\texttt{s} from a matrix via
\texttt{C.A = C.A[setdiff(1:end, s), :]} creates a \texttt{Set},
collects the complement indices, indexes the matrix, and allocates a new
matrix.  This is done inside the working-set update, which may execute
multiple times per iteration.

\paragraph{Fix.}
Replace with direct index construction:
\texttt{C.A = C.A[[1:s-1; s+1:end], :]}.
This avoids the \texttt{Set} overhead.  For vectors, replace with
\texttt{deleteat!} where the original is already a copy.

%-----------------------------------------------------------------------
\subsection{O5: Redundant allocations in \texttt{psi}}
\label{sec:o5}
%-----------------------------------------------------------------------

\juliaref{enlsip\_functions.jl}{1296--1331}

Each call to the merit function \texttt{psi} allocates
\texttt{x\_new = x + $\alpha$ * p} (a fresh $n$-vector),
\texttt{rx\_new = zeros(m)}, and \texttt{cx\_new = zeros(l)}.
During the line search, \texttt{psi} is called 3--10~times per
iteration, yielding 9--30 allocations of $O(m+\ell+n)$ vectors.

\paragraph{Fix.}
Add pre-allocated workspace vectors to \texttt{psi} as keyword arguments
with default allocation, so that callers in the hot path
(\texttt{linesearch\_constrained}, \texttt{goldstein\_armijo\_step}) can
pass reusable buffers.

%-----------------------------------------------------------------------
\subsection{O6: Redundant allocations in \texttt{linesearch\_constrained}
  and \texttt{compute\_steplength}}
\label{sec:o6}
%-----------------------------------------------------------------------

\juliaref{enlsip\_functions.jl}{1990, 2258--2259}

Several hot-path functions allocate fresh \texttt{zeros(T, m)} and
\texttt{zeros(T, l)} vectors that could be pre-allocated once outside
the main loop and passed through.  The functions
\texttt{linesearch\_constrained} and \texttt{compute\_steplength} each
allocate \texttt{rx\_new}, \texttt{cx\_new} vectors that are only used
locally.

\paragraph{Fix.}
Pre-allocate these vectors in the main \texttt{enlsip} function and
pass them as arguments.  This eliminates per-iteration allocations.

%-----------------------------------------------------------------------
\subsection{O7: Redundant \texttt{p\_gn = zeros(T, n)} per iteration}
\label{sec:o7}
%-----------------------------------------------------------------------

\juliaref{enlsip\_functions.jl}{2657, 2728}

The search direction buffer \texttt{p\_gn} is allocated as a fresh
zero vector at the start of every iteration.  Since it is immediately
overwritten by \texttt{gn\_search\_direction}, a single buffer allocated
once and filled in-place suffices.

\paragraph{Fix.}
Allocate once before the main loop; fill with zeros at each iteration
start using \texttt{fill!}.

%-----------------------------------------------------------------------
\subsection{O8: \texttt{v[1:m] = rx[:]} unnecessary copy in \texttt{concatenate!}}
\label{sec:o8}
%-----------------------------------------------------------------------

\juliaref{enlsip\_functions.jl}{1637}

\texttt{rx[:]} creates a copy of \texttt{rx}; since we are assigning
into a pre-existing slice of~\texttt{v}, the intermediate copy is
wasteful.  Replace with \texttt{copyto!} or direct view assignment.

\paragraph{Fix.}
\texttt{v[1:m] .= rx}.

%-----------------------------------------------------------------------
\subsection{O9: Redundant QR factorisation after working-set rebuild}
\label{sec:o9}
%-----------------------------------------------------------------------

\juliaref{enlsip\_functions.jl}{712--716, 731--735, 757--762}

When a constraint deletion is reverted (infeasible direction test),
the active-constraint matrix \texttt{C.A} is rebuilt from scratch and
a fresh QR factorisation is computed.  In the no-deletion branch
(\texttt{s == 0}), the original \texttt{F\_A} could be reused directly
instead of re-computing the QR.  However, since the second-order
multiplier estimate or subsequent deletion may modify the working set,
the factorisation is generally needed.

\medskip\noindent
\textit{Not modified --- the saving is marginal and the logic would
become fragile.}

%-----------------------------------------------------------------------
\subsection{Summary of optimisations}
%-----------------------------------------------------------------------

\begin{table}[ht]
\centering
\caption{Summary of performance optimisations.}
\label{tab:optim}
\small
\begin{tabular}{@{}lllp{6.5cm}@{}}
\toprule
\textbf{ID} & \textbf{File} & \textbf{Lines} & \textbf{Status} \\
\midrule
O1 & \texttt{cnls\_model.jl}       & 11--35  & \textbf{Applied}: parametric function types \\
O2 & \texttt{cnls\_model.jl}       & 155--174 & Documented only (API stability) \\
O3 & \texttt{enlsip\_functions.jl} & various & \textbf{Applied}: \texttt{maximum(abs, $\cdot$)} \\
O4 & \texttt{enlsip\_functions.jl} & various & \textbf{Applied}: direct index construction \\
O5 & \texttt{enlsip\_functions.jl} & 1296    & \textbf{Applied}: pre-allocated workspace \\
O6 & \texttt{enlsip\_functions.jl} & various & \textbf{Applied}: lifted allocations \\
O7 & \texttt{enlsip\_functions.jl} & 2657, 2728 & \textbf{Applied}: single buffer \\
O8 & \texttt{enlsip\_functions.jl} & 1637    & \textbf{Applied}: \texttt{.=} broadcast \\
O9 & \texttt{enlsip\_functions.jl} & various & Documented only \\
\bottomrule
\end{tabular}
\end{table}

%=======================================================================
\section{Sparse Matrix Support}
\label{sec:sparse}
%=======================================================================

For problems with structured (sparse) Jacobians, the current
all-dense storage is a significant bottleneck.  Consider the
Chained Rosenbrock test problem with $n=1000$ parameters, $m=1998$
residuals, and $l=998$ equality constraints.  The residual Jacobian
$J\in\mathbb{R}^{1998\times 1000}$ is tridiagonal ($\sim$\,4\,000
non-zeros vs.\ 2\,000\,000 entries in the dense matrix), and the
constraint Jacobian $A\in\mathbb{R}^{998\times 1000}$ is similarly
banded.  Dense storage costs $\sim$\,24\,MB per matrix pair;
sparse storage costs $\sim$\,300\,KB.  Matrix--vector products
($J\,p$, $J^\top r$, $A\,p$) dominate the per-iteration cost and
scale as $O(mn)$ dense versus $O(\mathrm{nnz})$ sparse.

The changes described below allow user-provided Jacobians in any
\texttt{AbstractMatrix} format (including
\texttt{SparseMatrixCSC}) to flow through the solver without
unnecessary densification, while keeping the internal
active-constraint subsystem dense where required by the QR
factorisation.

%-----------------------------------------------------------------------
\subsection{S1: \texttt{AbstractMatrix} type signatures}
%-----------------------------------------------------------------------

All function signatures that accept the residual Jacobian $J$ or the
full constraint Jacobian $A$ as \emph{read-only} inputs (i.e.\
they only perform matrix--vector products or column/row reads) are
widened from \verb|Matrix{T}| to \verb|AbstractMatrix{T}|.  This
single change is sufficient for Julia's multiple dispatch to
select the appropriate sparse BLAS kernels for products like
\verb|J * p| or \verb|transpose(J) * rx|.

Affected functions include \texttt{new\_point!},
\texttt{gn\_search\_direction}, \texttt{sub\_search\_direction},
\texttt{newton\_search\_direction},
\texttt{second\_lagrange\_mult\_estimate!},
\texttt{search\_direction\_analys},
\texttt{update\_working\_set},
\texttt{compute\_steplength}, and
\texttt{upper\_bound\_steplength}.

%-----------------------------------------------------------------------
\subsection{S2: Return-value Jacobian evaluation}
%-----------------------------------------------------------------------

The original evaluation wrappers pre-allocate dense buffers and
copy into them:
\begin{lstlisting}
J = zeros(T, m, n)
J[:] = r.jacres_eval(x)
\end{lstlisting}
This forces densification of any sparse result.
The revised pattern removes the pre-allocation and uses the
matrix returned by the user function directly:
\begin{lstlisting}
J = r.jacres_eval(x)
\end{lstlisting}
Since the user function already allocates its return value,
eliminating the copy is more efficient even for dense
Jacobians.  The evaluation wrappers
\texttt{jacres\_eval} and \texttt{jaccons\_eval} now return the
matrix instead of writing in-place, while
\texttt{res\_eval!} and \texttt{cons\_eval!} remain in-place
(vectors are always dense).

%-----------------------------------------------------------------------
\subsection{S3: Sparse box-constraint Jacobians}
%-----------------------------------------------------------------------

The \texttt{box\_constraints} helper constructs Jacobians of the
form $\pm I_n$ restricted to rows with finite bounds.  The
original code builds a full dense $n\times n$ identity matrix and
then selects rows:
\begin{lstlisting}
Matrix{T}(I, n, n)[finite_rows, :]
\end{lstlisting}
For $n=1000$, this allocates an 8\,MB dense matrix only to
extract a handful of rows.  The replacement uses
\texttt{sparse}:
\begin{lstlisting}
sparse(I, n, n)[finite_rows, :]
\end{lstlisting}
The resulting Jacobian is a \texttt{SparseMatrixCSC} with at most
$n$ non-zeros, regardless of which bounds are finite.  Since these
Jacobians are constant (independent of $x$), they could also be
pre-computed once; the current implementation re-evaluates them on
each call but the sparse construction is already $O(n)$.

%-----------------------------------------------------------------------
\subsection{S4: Dense active-constraint interface}
%-----------------------------------------------------------------------

The active-constraint Jacobian \verb|C.A| (a $t\times n$ submatrix of $A$,
where $t$ is the current working-set size) remains a dense \verb|Matrix{T}|.
This is necessary because:
\begin{itemize}[nosep]
  \item The rank-revealing QR factorisation
    \verb|qr(C.A', ColumnNorm())| requires a dense argument in
    Julia's standard library.
  \item Row deletion, insertion, and diagonal scaling are performed
    frequently and are more efficient on dense storage.
  \item The matrix is small ($t \le n$, often $t \ll l$).
\end{itemize}
When extracting \verb|C.A| from the (potentially sparse) full
Jacobian $A$, an explicit conversion is performed:
\begin{lstlisting}
C.A = Matrix(A[active_indices, :])
\end{lstlisting}
This ensures that all downstream QR and triangular-solve
operations receive dense inputs.

%-----------------------------------------------------------------------
\subsection{S5: Sparse-aware Hessian finite differences}
%-----------------------------------------------------------------------

The Hessian approximation routines \texttt{hessian\_res!} and
\texttt{hessian\_cons!} accept $J$ or $A$ only indirectly (they
call the evaluation functions).  Since the Hessian $B$ itself is
dense (the finite-difference stencil generally fills all entries),
no sparse treatment is needed for the output.  However, the
function signatures are widened to accept sparse inputs where
they appear.

%-----------------------------------------------------------------------
\subsection{Summary of sparse changes}
%-----------------------------------------------------------------------

\begin{table}[ht]
\centering
\caption{Summary of sparse matrix support changes.}
\label{tab:sparse}
\small
\begin{tabular}{@{}lllp{6.5cm}@{}}
\toprule
\textbf{ID} & \textbf{File} & \textbf{Scope} & \textbf{Description} \\
\midrule
S1 & \texttt{enlsip\_functions.jl} & signatures & \texttt{Matrix\{T\}} $\to$ \texttt{AbstractMatrix\{T\}} \\
S2 & \texttt{cnls\_model.jl}       & eval.\ pattern & return-value Jacobian evaluation \\
S3 & \texttt{cnls\_model.jl}       & box constraints & \texttt{sparse(I,n,n)} instead of \texttt{Matrix(I,n,n)} \\
S4 & \texttt{enlsip\_functions.jl} & active set & \texttt{Matrix(A[idx,:])} dense conversion \\
S5 & \texttt{enlsip\_functions.jl} & Hessians & widened input signatures \\
\bottomrule
\end{tabular}
\end{table}

\end{document}
